\scanpage{187}
\uline{Proposition} \quad Soit $f:X\unsure{\mathbin{-}|\to}Y$ un morphisme, avec

\begin{tabular}{|p{.86\linewidth}}
1) \deleted{$f$ loc de présentation finie et \unsure{loc} quasi-fini} \inserted{$Y$ loc. noeth et \unsure{univ.} japonais}\\[.25em]
2) $\deleted{\unsure{f}}X$ \deleted{irréductible} n'a qu'un nb fini de comp. irréd.\\[.25em]
3) $f$ loc. \inserted{de type fini,} quasi-fini \deleted{et \unsure{sé}}, $f$ \uline{séparé}
\end{tabular}

\begin{tabular}{|p{.86\linewidth}}
Sous ces conditions, $f$ est \uline{de type fini}
\end{tabular}

$-$

On peut supposer $\deleted{\unsure{S|X}}Y$ affine donc noeth, \deleted{\unsure{$f$}} $X$ intègre, \deleted{et $f$ dominant} \inserted{(car il suffit de voir que $f$ est qu. \unsure{cpt|cqt})}. Soit $U$ un ouvert \inserted{affine} non vide de $X$, il est quasi-fini séparé sur $Y$, donc par le main theorem il est $Y$-isomorphe
\diagram{\begin{tikzcd}[ampersand replacement=\&,column sep=1.25em,row sep=.55em]
\& X''\arrow[dl,"\text{fini}"']\arrow[dd]\\
X\arrow[dd]\&\\
\& X'\arrow[dl,"\substack{\text{fini}\\\text{\unsure{surj}}}"]\\
Y\&
\end{tikzcd}}
à un ouvert d'un \deleted{\unsure{fini $X'$ de $Y$}} \inserted{$Y$-}préschéma \uline{fini} $X'$, qu'on peut supposer intègre. Soit $X''$ \unsure{le} sous-préschéma fermé intègre de $X\times_Y X'$ \deleted{\unsure{propre de} \illegible{1}} \inserted{adhérence de} la ``diagonale'' de $U\times_Y U'$, donc $X''$ s'envoie dans $X$ et $X'$. Alors $X''\to X$ est fini, et étant dominant, est \uline{surjectif}. Il suffit donc de voir que $X''$ est de t.f. sur $Y$, ou \unsure{même} sur $X'$. Or $X''\to X'$ est un morphisme birationnel, \unsure{avec|car} $X''$ et $X'$ intègres. Donc \deleted{\unsure{$X'$|il}}il y a un $X'$-morphisme \uline{surjectif} $\widetilde X'\to X''$, où $\widetilde X'$ est le normalisé. Or $\widetilde X'$ est fini sur $Y$, on gagne .....

% archive sheet number: 186
